\documentclass[journal]{IEEEtran}

\usepackage{cite}
\usepackage{amsmath,amssymb,amsfonts}
\usepackage{graphicx}
\usepackage{textcomp}
\usepackage{xcolor}
\usepackage{listings}
\usepackage[inline, shortlabels]{enumitem}

\lstset{
	basicstyle=\footnotesize\ttfamily,
	breaklines=true,
	frame=single,
	captionpos=b,
	tabsize=2,
	showstringspaces=false,
	commentstyle=\color{gray},
	keywordstyle=\color{blue},
	stringstyle=\color{red},
	frame=none
}

\begin{document}
	
	\title{Demonstrating Good Practice in Coding}
	
	\author{Mindaugas Taujanskas}
	
	\maketitle
	
	\section{Don't Repeat Yourself (DRY)}
	DRY is a simple practice which promotes refactoring commonly used procedures into functional units which can be reused either to minimise verbatim repetition, or to parameterise a certain routine into a more generic building block. Consider an example from my removals system:
	
	\begin{lstlisting}[language=Python]
class ValidationMixin:
	def is_optional(self) -> bool:
		self._is_optional = getattr(
			self, "_is_optional", False
		)
		return self._is_optional
	
	def set_optional(self, to: bool) -> None:
		self._is_optional = to
		return self
	
	def set_validation_trigger(self, fn) -> None:
		self._validation_trigger = fn
		return self
	
	def register_validation_func(self, fn: ValidationFnT) -> None:
		if not hasattr(self, "state"):
			raise RuntimeError("Base class does not inherit StyledWidget")
		if not hasattr(self, "_validation_trigger"):
			raise RuntimeError(
				"Tried to register validation function with no trigger"
			)
		self._validation_fn = fn
		self._validation_trigger.connect(self._validate_callback)
		return self
	
	def _validate_callback(self) -> None:
		data = self.serialize()
		if not data and self.is_optional():
			self.state.emit("")
			return
		self.state.emit(
			"" if self._validation_fn(data) else "error"
		)
	
	def is_valid(self) -> bool:
		data = self.serialize()
		if not data:
			return self.is_optional()
		if not hasattr(self, "_validation_fn"):
			return True
		return self._validation_fn(data)
	
	def serialize(self) -> object:
		raise NotImplementedError
	\end{lstlisting}
	
	My rationale behind creating the \lstinline|ValidationMixin| class is due to earlier stages of development where I found myself recreating methods to validate form data (like emails, passwords, numbers,) and this became a very quick cause for concern when I found myself creating improved iterations of essentially the same validation routines over time, but leaving older iterations remnant, thus meaning I'd have to go back and rewrite them.
	As the number of UI field components grew it created increasingly more technical debt, so I extracted the key methods and logic into a mix-in class (technically-speaking, classes that \textit{specifically} extend the base functionality of another class), and subclassed it against every field component so that their behaviour would be guaranteed to be consistent.
	
	\subsection{Separation of concerns}
	
	A certain aspect where DRY may seem to be (falsely) violated, is with the model-view-controller (MVC) pattern. Consider the login procedure implemented at three stages, firstly as a SQL procedure:
	
	\begin{lstlisting}[language=SQL]
CREATE FUNCTION login_user(p_email TEXT, p_password TEXT)
BEGIN
	normalized_email = normalize_email(p_email);
	
	SELECT user_id, password_hash, user_role
	INTO stored_user_id, stored_hash, stored_user_role
	FROM Users
	WHERE email = normalized_email;
	
	IF NOT FOUND THEN
		RETURN QUERY
		SELECT 'Invalid email or password', NULL, NULL;
	END IF;
	
	IF crypt(p_password, stored_hash) = stored_hash THEN
		RETURN QUERY
		SELECT
			'',
			sign(/* ... */),
			stored_user_role;
	END IF;
	
	RETURN QUERY
	SELECT 'Invalid email or password', NULL, NULL;
END;
	\end{lstlisting}
	
	Which validates the email parameter, verifies it exists, and then validates the password and returns a session context once everything seems OK. Then the model implementation for the database in Python:
	
	\begin{lstlisting}[language=Python]
def proc_login_user(email: str, password: str) -> list[DictRow]:
	return call_proc("login_user", (email, password))
	\end{lstlisting}
	
	This closes together the boundary between database and application layer. Finally, in the user model:
	\newpage
	\begin{lstlisting}[language=Python]
class User:
	def __init__(
		self,
		email: str,
		password: str,
		*,
		token_role_pair: tuple[str, str] | None = None
	) -> None:
		if token_role_pair is not None:
			self.token, self.role = token_role_pair
			return
		error, *token_role_pair = db.proc_login_user(email, password)
		self.maybe_raise_exception(error)
		self.token, self.role = token_role_pair
	\end{lstlisting}
	
	Could we not have just used \lstinline|call_proc| ourselves and saved the need to create a separate module? Better yet, we could've just concocted our own SQL \lstinline[language=SQL]|INSERT| statement in the \lstinline|User| class, and saved ourselves from nesting two layers of interaction.
	
	\textbf{No}. Each layer adds another level of flexibility and separates responsibility for:
	\begin{enumerate*}[(a)]
		\item who must maintain database-level security
		\item how the database interface is implemented (PostgreSQL, MSSQL?)
		\item what it means to login, and what the fail-states are.
	\end{enumerate*}
	
	DRY takes least precedence over overarching principles such as the separation of concerns, it tends to follow as a result of being disciplined to other good practices.
	
	\section{Keep It Simple, Stupid (KISS)}
	KISS retains its relevancy at all stages of development. It is without a doubt that software projects will become increasingly complex as they grow, but ideally only as the sum of their parts, and not because of obscure, cryptic parts which form an ugly hornet's nest of complex interactions and couplings.
	
	A small snippet of my project which I found a little averse to this principle is the following:
	
	\begin{lstlisting}[language=Python]
def get_widget(self, name: str, default: None = None) -> "QWidget | None":
	# NOTE: :)
	if any((which := field).name == name for field in self.fields):
		return which
	return default
	\end{lstlisting}
	
	The canonical form for writing this would be simply:
	
	\begin{lstlisting}[language=Python]
def get_widget(self, name: str, default: None = None) -> "QWidget | None":
	for field in self.fields:
		if field.name == name:
			return field
	return default
	\end{lstlisting}
	
	But the former read left-to-right explains succinctly the mechanism of the function, let alone the function's name itself, although it utilises a fairly cryptic feature of the Python language.
	
	\section{The SOLID Principles}
	SOLID is a very well-studied and understood set of principles that underline object-oriented design, which outlines how classes should be created, maintained, constrained, related and designed. It promotes loosely-coupled classes which subsequently weighs positively in on their maintainability, modularity and testing.
	
	I won't go into detail on the specific principles, but they are a quintessential baseline for developing a well-oiled object-oriented application.
\end{document}